% TEMPLATE-MILC-v3.tex - plantilla maestra UNICA de las guias MILC v3.
%
% La usan las cuatro familias de guias, cada una con su driver:
%   · Grados 6-11          -> scripts/build-guias-g11.py
%   · Bebras               -> scripts/build-guias-bebras.py
%   · Semillero            -> scripts/build-guias-semillero.py
%   · Territorio interior  -> scripts/build-guias-territorio-interior.py
%
% Antes habia tres copias casi identicas (~400 lineas iguales, 6 distintas):
% cada mejora habia que aplicarla tres veces, y en la practica no se hacia
% (el motor de recursos vivia solo en dos de ellas). Lo que varia entre
% familias esta parametrizado en 6 placeholders que cada driver llena
% (se nombran aqui SIN su sintaxis real: el builder reemplaza por texto plano,
% tambien dentro de los comentarios, y un valor multilinea rompe el preambulo):
%
%   HEADER_IZQ        encabezado izquierdo de pagina
%   HEADER_DER        encabezado derecho de pagina
%   PORTADA_ETIQUETA  rotulo sobre el numero grande (GRADO/MOMENTO/MODULO)
%   PORTADA_SERIE     linea de serie de la portada
%   PORTADA_ID        identificador de la guia en la portada
%   PORTADA_CIRCULO   el circulito de grado (°), o vacio si no aplica
%   PDF_TITULO        titulo en texto plano (sin \textbf ni comillas TeX) para
%                     los metadatos del PDF (/Title); lo produce lib_xelatex.texto_pdf
%
% Composicion (auditoria 2026-09-03): las bandas de estacion piden espacio
% (\Needspace*) y prohiben el salto justo despues (\nopagebreak), asi no quedan
% huerfanas al pie; las cajas largas (softbox, drawbox, infoband, puentebox) son
% `breakable`, asi una caja de media pagina ya no arrastra todo a la siguiente
% dejando la anterior al 40 %. Todo texto sobre mostaza o verde va en milcNegro
% (blanco sobre #E5B400 da 1,93:1 y sobre #58A923 2,95:1; ninguno pasa WCAG AA).
% El resto de claves las documenta content/guias/_SCHEMA.md. El script valida
% que no queden placeholders sin reemplazar antes de escribir el archivo final.

% Etiquetado PDF (latex-lab): /Lang, estructura y texto alternativo de las
% figuras, para lectores de pantalla. Debe ir ANTES de \documentclass.
\DocumentMetadata{lang=es-CO,tagging=on}
\documentclass[11pt,letterpaper]{article}
% headheight=24pt: la cabecera derecha lleva dos lineas (identidad de la guia
% y estacion actual). headsep se acorta en la misma medida para que el bloque
% de texto quede exactamente donde estaba con headheight=12pt.
\usepackage[margin=2.0cm,headheight=24pt,headsep=13pt]{geometry}
\usepackage[table]{xcolor}
\usepackage{fontspec}
\usepackage[spanish]{babel}
\usepackage{tikz}
% Librerías TikZ para los diagramas de las guías (bloque `recursos:`):
%  · babel  → babel en español vuelve activos caracteres como «>», y TikZ los usa
%             en sus opciones (>=stealth, ->). Sin esta librería, cualquier
%             diagrama con flechas revienta con «\language@active@arg> has an
%             extra }». Debe ir ANTES de que se dibuje nada.
%  · positioning        → sintaxis «below left=2pt and 3pt».
%  · decorations.pathreplacing → llaves ({brace}) para señalar rangos.
\usetikzlibrary{babel,positioning,decorations.pathreplacing}
\usepackage{graphicx}
% Circuitos: el trabajo de electrónica se hace hoy con micro:bit y sus sensores
% integrados (MakeCode, sin cableado), así que no se necesitan esquemáticos.
% Si algún día entran montajes con protoboard, basta descomentar esta línea:
% los diagramas .tex/.tikz declarados en `recursos:` ya se incrustan solos.
% \usepackage{circuitikz}
\usepackage[most]{tcolorbox}
\usepackage{tabularx}
\usepackage{array}
\usepackage{fancyhdr}
\usepackage{titlesec}
\usepackage[document]{ragged2e}  % bandera a la derecha en todo el cuerpo
\usepackage{enumitem}
\usepackage{microtype}
\usepackage{etoolbox}
\usepackage{needspace}
% hyperref al final del preambulo: metadatos del PDF (titulo, autor, idioma,
% familia), marcadores por estacion (\pdfbookmark) y enlaces sin recuadro.
\usepackage[hidelinks,
  pdftitle={Referencias relativas y absolutas — fija la rejilla, repite la puntada},
  pdfauthor={PhD. Álvaro Cárdenas Orozco},
  pdfsubject={Tecnología e Informática · Grado 8 · Guía 5 · MILC v3},
  pdflang={es-CO},
  bookmarksopen=true,bookmarksnumbered=false]{hyperref}

% Helvetica Neue no trae la flecha U+2192, asi que cada «→» escrito literalmente
% en el YAML se compilaba a NADA, en silencio (462 flechas en 61 guias: rutas del
% tipo «paleta → tipografia → lockup» salian rotas). Mapeamos el caracter a su
% version matematica, que si tiene glifo. Asi el autor puede seguir escribiendo
% «→» comodamente en el YAML.
\usepackage{newunicodechar}
\newunicodechar{→}{\ensuremath{\rightarrow}}
% Mismo problema con los simbolos de marca y vinetas: el «✓» de «una palomita ✓»
% salia como cuadrito vacio en el PDF de todas las guias que lo usaban.
\usepackage{pifont}
\newunicodechar{✓}{\ding{51}}
\newunicodechar{✗}{\ding{55}}
\newunicodechar{✔}{\ding{52}}
\newunicodechar{•}{\textbullet}
\newunicodechar{≥}{\ensuremath{\geq}}
\newunicodechar{≤}{\ensuremath{\leq}}

\setmainfont{Helvetica Neue}
\setsansfont{Helvetica Neue}
\setlength{\parindent}{0pt}
\setlength{\parskip}{6pt}
% Sin particion de palabras: en 6.º-7.º una palabra cortada por guion al final
% de linea («Comuni-cacion») frena la lectura mucho mas que una linea corta.
\hyphenpenalty=10000
\exhyphenpenalty=10000
\pretolerance=2000
\tolerance=3000
\setlength{\emergencystretch}{3em}
\linespread{1.10}
\renewcommand{\arraystretch}{1.25}
\setlist{leftmargin=5mm,itemsep=1mm,topsep=1mm}
\newcolumntype{Y}{>{\RaggedRight\arraybackslash}X}

\definecolor{milcMagenta}{HTML}{E6007E}
\definecolor{milcVino}{HTML}{5A0038}
\definecolor{milcTurquesa}{HTML}{0093A5}
\definecolor{milcVerde}{HTML}{58A923}
\definecolor{milcMostaza}{HTML}{E5B400}
\definecolor{milcOcre}{HTML}{B86600}
\definecolor{milcNegro}{HTML}{241816}
\definecolor{milcGris}{HTML}{F7F7F4}
\definecolor{milcLinea}{HTML}{E8E2DD}
\definecolor{milcGradePrimary}{HTML}{FF6600}
\definecolor{milcGradeDark}{HTML}{9A3E00}
\definecolor{milcGradeSoft}{HTML}{FFE4D0}
\definecolor{milcGradeAccent}{HTML}{CFE7FF}
\definecolor{milcGradeText}{HTML}{FFFFFF}
\color{milcNegro}

\pagestyle{fancy}
\fancyhf{}
% Cabecera de dos lineas: identidad de la guia arriba y, a la derecha y en
% gris, la estacion en curso (\markright desde \milcestacion). Antes de la
% primera estacion la segunda linea va vacia; el \strut mantiene la altura
% para que la linea de arriba no baile entre paginas.
\lhead{\small Tecnología e Informática · Grado 8\\[-1pt]{\footnotesize\strut}}
\rhead{\small Guía 5 · MILC v3\\[-1pt]{\footnotesize\strut\color{milcNegro!65}\rightmark}}
\cfoot{\small\thepage}
\renewcommand{\headrulewidth}{0pt}

% Sobre mostaza (#E5B400) y verde (#58A923) el texto va en milcNegro; sobre el
% resto de la paleta, en blanco. Se usa dentro de fonttitle/fontupper, donde un
% \color posterior manda sobre coltitle.
\newcommand{\milctintasobre}[1]{%
  \ifboolexpr{test{\ifstrequal{#1}{milcMostaza}} or test{\ifstrequal{#1}{milcVerde}}}%
    {\color{milcNegro}}{\color{white}}}

\newtcolorbox{titlebox}[1]{
  enhanced,colback=#1,colframe=#1,arc=7mm,boxrule=0pt,
  left=7mm,right=7mm,top=3mm,bottom=3mm,
  before skip=8pt,after skip=8pt,
  fontupper=\sffamily\bfseries\Large\color{white}
}

\newtcolorbox{softbox}[3]{
  enhanced,breakable,pad at break=3mm,
  colback=#2,colframe=#1,borderline west={4pt}{0pt}{#1},
  arc=2mm,boxrule=.4pt,left=4mm,right=4mm,top=3mm,bottom=3mm,
  before skip=6pt,after skip=8pt,title={#3},coltitle=white,
  colbacktitle=#1,fonttitle=\sffamily\bfseries\milctintasobre{#1},fontupper=\RaggedRight,
  attach boxed title to top left={xshift=3mm,yshift=-2mm},
  boxed title style={arc=2mm, boxrule=0pt}
}

\newtcolorbox{drawbox}[2]{
  enhanced,breakable,pad at break=4mm,
  colback=white,colframe=#1,arc=4mm,boxrule=1pt,
  left=5mm,right=5mm,top=4mm,bottom=4mm,before skip=8pt,after skip=8pt,
  title={#2},coltitle=white,colbacktitle=#1,fonttitle=\sffamily\bfseries\milctintasobre{#1},
  fontupper=\RaggedRight,
  attach boxed title to top left={xshift=4mm,yshift=-2mm},
  boxed title style={arc=3mm, boxrule=0pt}
}

\newtcolorbox{infoband}[1]{
  enhanced,breakable,pad at break=2mm,colback=#1!8,colframe=#1!50,arc=2mm,boxrule=.5pt,
  left=4mm,right=4mm,top=2mm,bottom=2mm,before skip=4pt,after skip=4pt,
  fontupper=\small\RaggedRight
}

% Puente narrativo entre secciones (contrato editorial Conéctate).
% Da continuidad: "Acabas de X. Ahora vas a Y porque Z."
% Visualmente: banda izquierda en color del grado, texto en itálica pequeña.
\newtcolorbox{puentebox}{
  enhanced,breakable,pad at break=2mm,colback=milcGradeSoft,colframe=milcGradeDark,
  borderline west={3pt}{0pt}{milcGradeDark},
  arc=0mm,boxrule=0pt,left=5mm,right=4mm,top=2.5mm,bottom=2.5mm,
  before skip=4pt,after skip=8pt,
  fontupper=\itshape\small\color{milcNegro}\RaggedRight
}
% La flecha va en modo matematico a proposito: Helvetica Neue NO tiene el glifo
% U+2192, asi que \textrightarrow se compilaba a NADA («Missing character: There
% is no → in font Helvetica Neue») y el puente salia sin flecha. En modo
% matematico la flecha viene de la fuente matematica y si se dibuja.
\newcommand{\puente}[1]{%
  \begin{puentebox}%
    \textcolor{milcGradeDark}{$\rightarrow$}\hspace{2mm}#1%
  \end{puentebox}%
}

\newcommand{\linead}[1]{\noindent\color{milcLinea}\rule{#1}{0.5pt}\color{milcNegro}}
\newcommand{\respuesta}[1]{\par\vspace{2mm}\foreach \n in {1,...,#1}{\noindent\color{milcLinea}\rule{\linewidth}{0.5pt}\par\vspace{5mm}}\color{milcNegro}}
\newcommand{\checkbox}{\textcolor{milcGradeDark}{\fbox{\phantom{x}}}\hspace{5pt}}

% ─── Listas aireadas para las guías socioemocionales ────────────────────
% Componer las verificaciones como «\checkbox texto\\» deja el texto que salta
% de línea debajo del cuadrito y apelmaza el bloque. Estos entornos alinean la
% continuación y airean los ítems. Los usa Territorio Interior; cualquier
% familia puede hacerlo.
\newlist{checklista}{itemize}{1}
\setlist[checklista]{
  label=\textcolor{milcGradeDark}{\fbox{\phantom{x}}},
  leftmargin=9mm, labelsep=3.2mm, itemsep=2.4mm,
  topsep=2.5mm, parsep=0pt, partopsep=0pt, before=\RaggedRight
}
\newlist{pasoslista}{enumerate}{1}
\setlist[pasoslista]{
  label=\textcolor{milcGradeDark}{\sffamily\bfseries\arabic*.},
  leftmargin=8mm, labelsep=3mm, itemsep=2.8mm,
  topsep=1mm, parsep=0pt, partopsep=0pt, before=\RaggedRight
}

\newcommand{\phasecard}[4]{
\begin{tcolorbox}[enhanced,colback=#3,colframe=#2,arc=3mm,boxrule=.5pt,height=3.05cm,valign=center,halign=center,left=1.5mm,right=1.5mm,top=2mm,bottom=2mm]
{\sffamily\bfseries\fontsize{22}{24}\selectfont\textcolor{#2}{#1}}\par
{\sffamily\bfseries\small #4}
\end{tcolorbox}
}

% ── Piezas del rediseno 2026 (legibilidad 6.º-11.º) ───────────────────
% Banda de estacion: reemplaza el titlebox macizo. Titulo a la izquierda,
% dato practico (tiempo/modalidad) a la derecha, en una sola linea.
% Nunca huerfana: pide 9 lineas libres (o salta de pagina rellenando con
% \vfill) y prohibe el salto justo despues de la banda. Nueve y no seis:
% la caja partible que sigue necesita ~80pt (titulo adosado, relleno y dos
% lineas) para arrancar en la misma pagina; con menos, tcolorbox la manda
% entera a la siguiente y la banda se queda sola. El penalty 10000 va
% en `after`, ANTES del espacio posterior: un `after skip` (aunque sea 0pt)
% es un \vskip tras la caja, y ese glue es un punto de corte legal que un
% \nopagebreak escrito despues de \end{tcolorbox} ya no cierra. Cada banda
% deja marcador PDF y actualiza la cabecera derecha.
\newcounter{milcestacion}
\newcommand{\milcestacion}[2]{%
\Needspace*{9\baselineskip}%
\stepcounter{milcestacion}%
\pdfbookmark[1]{#2}{milcestacion\themilcestacion}%
\markright{#2}%
\begin{tcolorbox}[enhanced,colback=#1,colframe=#1,arc=2mm,boxrule=0pt,
  left=5mm,right=5mm,top=2.6mm,bottom=2.6mm,before skip=11pt,
  after={\par\nopagebreak\vskip7pt}]
{\sffamily\bfseries\fontsize{15}{18}\selectfont\milctintasobre{#1}#2}
\end{tcolorbox}}

% Tarjeta de paso numerada: sustituye las tablas «Pilar / Aplicacion», que
% eran formato de consulta docente, no de estudiante. El numero va en un
% circulo sobre el borde para que la secuencia se lea de un vistazo.
\newcommand{\milcpaso}[3]{%
\Needspace{4\baselineskip}%
\begin{tcolorbox}[enhanced,colback=white,colframe=milcLinea,arc=1.5mm,boxrule=.9pt,
  left=14mm,right=4mm,top=3mm,bottom=3mm,before skip=4pt,after skip=4pt,
  fontupper=\RaggedRight,
  overlay={\node[anchor=north west,circle,fill=#2,text=white,inner sep=0pt,
    minimum size=7.4mm,font=\sffamily\bfseries\small]
    at ([xshift=3.6mm,yshift=-3.2mm]frame.north west) {#1};}]
#3
\end{tcolorbox}}

% Casilla marcable de tamano util con lapiz (la anterior era \fbox{\phantom{x}},
% de alto variable segun el texto que la seguia).
\newcommand{\milcmarca}{\framebox[3.8mm]{\rule{0pt}{3.4mm}}}

% Fila marcable de lista suelta (checks de la estacion 1).
\newcommand{\milccheckrow}[1]{%
\par\vspace{1.8mm}\noindent
\makebox[6.4mm][l]{\raisebox{-.55mm}{\textcolor{milcNegro}{\milcmarca}}}%
\parbox[t]{\dimexpr\linewidth-6.4mm\relax}{\RaggedRight #1}\par}

% Celda marcable dentro de la rubrica (tabla de autoevaluacion). Misma
% sangria francesa, pero pensada para vivir dentro de una columna X.
\newcommand{\milcopcion}[1]{%
\makebox[5.2mm][l]{\raisebox{-.55mm}{\textcolor{milcNegro}{\milcmarca}}}%
\parbox[t]{\dimexpr\linewidth-5.2mm\relax}{\RaggedRight #1}}

% ── Recursos visuales de la guía ──────────────────────────────────────
% Declarados en el bloque `recursos:` del YAML y emitidos por el builder.
% \guiaFigura   → imagen rasterizada o PDF (foto, carta celeste, montaje).
% \guiaDiagrama → fuente TikZ/circuitikz (.tex) incrustada como vector:
%                 es la vía para circuitos y esquemáticos editables.
% [#1] = texto alternativo (alt) · #2 = ruta relativa al .tex · #3 = pie
% El alt llega al PDF etiquetado (/Alt) y lo lee el lector de pantalla.
\newcommand{\guiaFigura}[3][]{%
\begin{center}
\includegraphics[width=0.88\linewidth,height=0.38\textheight,keepaspectratio,alt={#1}]{#2}\\[1.5mm]
{\footnotesize\itshape\textcolor{milcNegro!75}{#3}}
\end{center}
}

\newcommand{\guiaDiagrama}[2]{%
\begin{center}
\input{#1}\\[1.5mm]
{\footnotesize\itshape\textcolor{milcNegro!75}{#2}}
\end{center}
}

\begin{document}
\thispagestyle{empty}

% ── Portada ───────────────────────────────────────────────────────────
% Antes: un rectangulo de color a pagina completa. Se veia bien en pantalla,
% pero en la fotocopiadora del colegio se comia el toner y salia gris sucio.
% Ahora la banda ocupa el tercio superior y el resto va en blanco.
\begin{tikzpicture}[remember picture,overlay]
  \path[fill=milcGradePrimary]
    (current page.north west) rectangle ([yshift=-10.6cm]current page.north east);

  \node[anchor=north west,text=milcGradeText] at ([xshift=2.0cm,yshift=-1.55cm]current page.north west)
    {\sffamily\bfseries\fontsize{12}{14}\selectfont GRADO};
  \node[anchor=north west,text=milcGradeText,opacity=.85] at ([xshift=2.0cm,yshift=-2.35cm]current page.north west)
    {\sffamily\bfseries\fontsize{8.5}{10}\selectfont SERIE GUÍAS · TECNOLOGÍA E INFORMÁTICA};
  \node[anchor=north east,text=milcGradeText,opacity=.85,align=right] at ([xshift=-2.0cm,yshift=-1.55cm]current page.north east)
    {\sffamily\bfseries\fontsize{9}{12}\selectfont I.E. Sor María Juliana\\Cartago, Valle del Cauca};

  \node[anchor=west,text=milcGradeText] (gradonum) at ([xshift=1.85cm,yshift=-7.1cm]current page.north west)
    {\sffamily\bfseries\fontsize{112}{112}\selectfont 8};
  % El circulito de grado (°) solo aplica a las guias de grado («11°»); en
  % Bebras y Semillero el numero grande es un momento/modulo, no un grado.
  % Se posiciona relativo al nodo `gradonum`, no en coordenadas absolutas.
  \draw[milcGradeText,line width=1.6pt]
    ([xshift=2.0mm,yshift=-4.5mm]gradonum.north east) circle (.15cm);

  \node[anchor=west,text=milcGradeText,text width=11.4cm,align=left] at ([xshift=1.5cm,yshift=0.5cm]gradonum.east)
    {
     {\sffamily\bfseries\fontsize{9}{11}\selectfont\MakeUppercase{Período 1 · Análisis de datos con phronesis}}\\[3mm]
     {\sffamily\bfseries\fontsize{21}{25}\selectfont\setlength{\baselineskip}{27pt}Fija la rejilla,\\[4pt]repite la puntada}\\[3.5mm]
     {\sffamily\bfseries\fontsize{15}{18}\selectfont Guía 5-8-TIC}
    };
\end{tikzpicture}

\vspace*{9.4cm}
\vfill

{\sffamily\bfseries\fontsize{8}{10}\selectfont\textcolor{milcGradeDark}{LO QUE TE LLEVAS HOY}}

\begin{tcolorbox}[enhanced,colback=white,colframe=milcNegro,arc=2mm,boxrule=1.6pt,
  left=5mm,right=5mm,top=4mm,bottom=4mm,before skip=3pt,after skip=12pt,
  fontupper=\RaggedRight\fontsize{11}{15.5}\selectfont]
Dos hojas de Excel. En la primera, el IVA y un descuento calculados sobre diez productos con una sola fórmula que se arrastra. En la segunda, la tabla de multiplicar del 1 al 10 construida con una única fórmula copiada a cien celdas. Y capturas que muestren que las fórmulas siguen bien al final del rango.
\end{tcolorbox}

\begin{tcolorbox}[enhanced,colback=milcGris,colframe=milcGris,arc=2mm,boxrule=0pt,
  left=5mm,right=5mm,top=3.5mm,bottom=3.5mm,after skip=12pt]
{\sffamily\bfseries\fontsize{8}{10}\selectfont\textcolor{milcNegro!55}{ESTA GUÍA ES DE}}\\[3mm]
\begin{tabularx}{\linewidth}{X p{3.2cm} p{3.2cm}}
\linead{\linewidth} & \linead{3.0cm} & \linead{3.0cm}\\[-1mm]
{\fontsize{8.5}{10}\selectfont\textcolor{milcNegro!55}{Nombre completo}} &
{\fontsize{8.5}{10}\selectfont\textcolor{milcNegro!55}{Grupo}} &
{\fontsize{8.5}{10}\selectfont\textcolor{milcNegro!55}{Fecha}}\\
\end{tabularx}
\end{tcolorbox}

\vfill

\noindent\textcolor{milcLinea}{\rule{\linewidth}{1pt}}\\[2mm]
{\sffamily\bfseries\fontsize{9.5}{12}\selectfont Autor: Dr. Álvaro Cárdenas Orozco}\\[1mm]
{\sffamily\fontsize{8.5}{11}\selectfont\textcolor{milcNegro!55}{Metodología MILC · apertura ancestral · pensamiento computacional · triángulo Dussel–estoicismo–Floridi}}

\newpage

\section*{Apertura: Referencias relativas y absolutas --- fija la rejilla, repite la puntada}

\begin{softbox}{milcGradeDark}{milcGradeSoft}{Saber ancestral}
En Ansermanuevo y Cartago hay un oficio que trabaja quitando en vez de poniendo: el \textbf{calado}. La caladora tensa la tela y retira hilos de la trama y de la urdimbre, de uno en uno. «Si tratas de deshilar dos hilos a la vez, ellos no te dejan, se pegan entre sí», cuenta Olivia, caladora de 65 años (Pérez-Bustos, 2019). Lo que queda no es un hueco: es una \textbf{rejilla regular}. Esa rejilla se fija primero y no se toca más. Solo entonces empieza la puntada, que se repite celda por celda, siempre igual, hasta llenar el diseño. Dos cosas, entonces: una que queda fija para toda la pieza, la rejilla, y otra que se copia y se mueve, la puntada. Si cambias el orden, la tela se arruina. La \textbf{cara de exclusión} del oficio: es un saber que se transmite en casa y sin título, y la caladora experta no aparece en ningún registro oficial. En Excel vas a hacer hoy lo mismo: fijar lo que no cambia con el signo \$, y dejar libre lo que sí se mueve al copiar.\par\smallskip{\footnotesize\textcolor{milcNegro!70}{\textbf{Fuente.} Cuéllar Barona, M., Sánchez-Aldana, E. y Pérez-Bustos, T. (Eds.). (2019). \emph{Papel de Colgadura}, 18. Universidad Icesi.}}
\end{softbox}

\begin{softbox}{milcTurquesa}{milcGris}{Saber contemporáneo}
Una \textbf{referencia} es la dirección de una celda dentro de una fórmula, como B2. Hay dos formas de escribirla. \textbf{Relativa}, B2: cuando copias la fórmula una fila abajo, se convierte en B3; la dirección se mueve con la copia. \textbf{Absoluta}, \$B\$2: cuando copias la fórmula, sigue apuntando a B2; el signo \$ la deja fija. Hay una tercera, la \textbf{mixta}: \$B2 fija la columna y deja libre la fila; B\$2 fija la fila y deja libre la columna. La regla del oficio es una sola: fija lo que no debe cambiar al copiar y deja libre lo que sí. Con eso, una fórmula bien escrita se copia a cien celdas sin error. Sin eso, escribes cien fórmulas a mano, y alguna sale mal.
\end{softbox}

\begin{softbox}{milcMostaza}{milcGris}{Pregunta-puente}
\textit{La caladora fija la rejilla una vez y después repite la misma puntada en cada celda. Cuando copies una fórmula a diez filas, ¿qué parte de ella es la rejilla, la que no debe moverse? ¿Y qué parte es la puntada, la que debe cambiar en cada fila?}
\end{softbox}

\begin{softbox}{milcVerde}{milcGris}{Saber y saber hacer}
Al terminar podrás: (1) \textbf{identificar} qué cambia y qué no cuando copias una fórmula con y sin \$; (2) \textbf{explicar} con tus palabras para qué sirve el signo \$; (3) \textbf{aplicar} referencias relativas, absolutas y mixtas en dos ejercicios reales; (4) \textbf{evaluar} una fórmula ajena y decir si se va a romper al copiarla.
\end{softbox}



\small
\begin{tabularx}{\textwidth}{>{\bfseries}p{3.0cm}Y}
\rowcolor{milcGradePrimary!12}Autor & Dr. Álvaro Cárdenas Orozco\\
DBA & Tratamiento de datos cuantitativos con criterio ético (MEN, T\&I)\\
Referentes & Phronesis aristotélica · Floridi (ética del dato) · Estoicismo (ver lo que es)\\
Producto final & Dos hojas de Excel. En la primera, el IVA y un descuento calculados sobre diez productos con una sola fórmula que se arrastra. En la segunda, la tabla de multiplicar del 1 al 10 construida con una única fórmula copiada a cien celdas. Y capturas que muestren que las fórmulas siguen bien al final del rango.\\
\end{tabularx}
\normalsize

\puente{En la sesión 4 escribiste tus primeras funciones. Hoy aprendes a copiarlas a muchas celdas sin que se rompan. En la sesión 6 vas a combinar operadores en una sola fórmula, y todo se apoya en las referencias de hoy.}

\milcestacion{milcGradeDark}{Ruta MILC: así vamos a aprender}
\begin{tabularx}{\textwidth}{>{\centering\arraybackslash}X>{\centering\arraybackslash}X>{\centering\arraybackslash}X>{\centering\arraybackslash}X}
\phasecard{1}{milcVerde}{milcGris}{Escucha\\[-1mm]\small Observo, escucho y pregunto.} &
\phasecard{2}{milcTurquesa}{milcGris}{Entender\\[-1mm]\small Sistematizo y ordeno.} &
\phasecard{3}{milcMagenta}{milcGris}{Hacer\\[-1mm]\small Construyo y aplico.} &
\phasecard{4}{milcVino}{milcGris}{Cierre\\[-1mm]\small Evalúo y reflexiono.}
\end{tabularx}


\puente{Antes de la teoría, vas a hacer un experimento de cuatro pasos en Excel y a ver con tus ojos qué le pasa a una fórmula cuando la arrastras, con \$ y sin \$.}

\milcestacion{milcVerde}{Estación 1 de 3 · Escucha}

\begin{softbox}{milcVerde}{milcGris}{Escena para escuchar}
\textbf{Actividad 1 · IDENTIFICA --- El experimento de fijar una celda} (15 min · individual). (1) En Excel escribe 10 en B2 y 20 en C2. En D2 escribe \texttt{=B2*C2}: debe dar 200. (2) Escribe 30 en B3 y 40 en C3. Arrastra la fórmula de D2 hasta D3: debe dar 1200. (3) En E2 escribe \texttt{=B2*\$C\$2}: da 200. Arrastra hasta E3: da 600, porque C2 quedó fija. (4) Haz clic en D3 y en E3 y mira la fórmula en la barra: ¿qué cambió en cada una? (5) Escribe en una línea la diferencia entre la columna D y la columna E.
\end{softbox}

{\sffamily\bfseries\fontsize{8}{10}\selectfont\textcolor{milcNegro!55}{MARCA CUANDO LO HAYAS HECHO}}
\milccheckrow{Hice los cuatro pasos y obtuve 200, 1200, 200 y 600.}
\milccheckrow{Miré las fórmulas de D3 y E3 en la barra y anoté qué cambió.}
\milccheckrow{Escribí en una línea qué hace el signo \$.}

\begin{infoband}{milcVerde}


\textbf{Lo que va al cuaderno (Actividad 1 · IDENTIFICA).} \textbf{Título:} «Actividad 1 --- El experimento de fijar una celda». \textbf{Formato:} ficha con los cuatro pasos, los cuatro resultados y una línea sobre la diferencia entre D y E. \textbf{Extensión:} media página. \textbf{Sabes que terminaste cuando} los cuatro resultados coinciden con los de la guía y puedes explicarle a un compañero qué hace el \$ con este ejemplo.
\end{infoband}

\puente{Ya viste el efecto del \$. Ahora entiendes las tres referencias, aprendes el atajo F4 y construyes con tu pareja la hoja del IVA y el descuento.}

\milcestacion{milcTurquesa}{Estación 2 de 3 · Entender}

\begin{softbox}{milcTurquesa}{milcGris}{Lo que vas a entender}
\textbf{Actividad 2 · APLICA --- IVA y descuento con una fórmula que se copia} (30 min · parejas). (1) Con tu pareja, escriban en la columna A diez productos de la tienda escolar y en la B sus precios. (2) En E1 escriban 19~\% y en E2 escriban 10~\%: son el IVA y el descuento. (3) En C2 escriban \texttt{=B2*\$E\$1} y arrastren hasta C11: es el IVA de cada producto. (4) En D2 escriban \texttt{=B2*(1-\$E\$2)} y arrastren hasta D11: es el precio con descuento. (5) Cambien el IVA de E1 a 21~\% y miren la columna C: si las diez filas cambiaron solas, la referencia quedó bien fijada.
\end{softbox}

\milcpaso{1}{milcTurquesa}{\textbf{Las tres referencias.} Relativa, B2: se mueve al copiar; sirve cuando cada fila tiene su propio dato, como el precio. Absoluta, \$B\$2: queda fija al copiar; sirve para constantes, como el IVA. Mixta, \$B2 o B\$2: fija solo la columna o solo la fila; sirve cuando copias en las dos direcciones.}
\milcpaso{2}{milcTurquesa}{\textbf{El atajo F4.} Pon el cursor sobre la referencia en la barra de fórmulas y presiona F4. Cada toque cambia entre B2, \$B\$2, B\$2 y \$B2. Es más rápido que escribir los signos a mano y evita olvidar uno.}
\milcpaso{3}{milcTurquesa}{\textbf{Dos preguntas antes de copiar.} Si arrastro hacia abajo, ¿esta dirección debe seguir en la misma fila o bajar con la copia? Si arrastro hacia la derecha, ¿debe seguir en la misma columna o moverse? Las respuestas dicen dónde va el \$.}
\milcpaso{4}{milcTurquesa}{\textbf{Verifica al final del rango.} La fórmula puede verse bien en la primera celda y estar rota en la última. Haz clic en la última celda del rango y mira la barra: si apunta a donde debe, la copia quedó bien.}

\begin{drawbox}{milcTurquesa}{Las tres referencias, en una mirada}
\textbf{Relativa} \texttt{B2}: se mueve al copiar. \\[1mm] \textbf{Absoluta} \texttt{\$B\$2}: queda fija. \\[1mm] \textbf{Mixta} \texttt{\$B2}: columna fija, fila libre. \\[1mm] \textbf{Mixta} \texttt{B\$2}: fila fija, columna libre. \\[1mm] \textbf{F4:} alterna entre las cuatro.
\end{drawbox}

\begin{softbox}{milcMostaza}{milcGris}{Errores comunes y su corrección}
(1) \textbf{Olvidar el \$ en la constante}: el IVA se mueve al copiar y apunta a celdas vacías. (2) \textbf{Fijar lo que debía moverse}: si fijas el precio, todas las filas calculan el mismo producto. (3) \textbf{Confundir la mixta}: fijar la fila cuando había que fijar la columna. (4) \textbf{No mirar la última celda}: el error se ve al final del rango, no al principio. (5) \textbf{Escribir cien fórmulas a mano}: alguna sale distinta y nadie la encuentra.
\end{softbox}

\begin{infoband}{milcTurquesa}


\textbf{Lo que va al cuaderno (Actividad 2 · APLICA).} \textbf{Título:} «Actividad 2 --- IVA y descuento con una fórmula». \textbf{Formato:} las dos fórmulas escritas tal cual, y una tabla de 2 filas y 2 columnas (qué fijamos con \$ / qué dejamos libre). \textbf{Extensión:} media página. \textbf{Sabes que terminaste cuando} al cambiar el IVA en E1 las diez filas de la columna C cambian solas, y la tabla dice qué se fijó y por qué.
\end{infoband}

\puente{Ya sabes fijar una constante. Ahora el reto grande: una tabla de multiplicar de cien celdas con una sola fórmula, y la revisión de la fórmula de tu pareja.}

\milcestacion{milcMagenta}{Estación 3 de 3 · Hacer}

\begin{softbox}{milcMagenta}{milcGris}{Cómo vas a construir tu producto}
\textbf{Actividad 3 · EVALÚA --- Cien celdas con una sola fórmula, y la fórmula de tu pareja} (25 min · individual). (1) En una hoja nueva escribe del 1 al 10 en B1 hasta K1, y del 1 al 10 en A2 hasta A11. (2) En B2 escribe \texttt{=\$A2*B\$1}. (3) Arrastra B2 hacia la derecha hasta K2. Después selecciona B2:K2 y arrastra hacia abajo hasta la fila 11. (4) Haz clic en K11 y mira la barra: debe decir \texttt{=\$A11*K\$1} y mostrar 100. (5) Intercambia la hoja con tu pareja y revisa su fórmula de B2: ¿se rompería al copiarla? Escribe tu veredicto en una línea.
\end{softbox}

\milcpaso{1}{milcMagenta}{\textbf{Tu entrega tiene tres piezas.} (1) La hoja del IVA y el descuento con las dos fórmulas copiadas. (2) La tabla de multiplicar con una única fórmula en cien celdas. (3) Las capturas de la última celda de cada rango y tu veredicto sobre la fórmula de tu pareja.}
\milcpaso{2}{milcMagenta}{\textbf{Por qué \$A2*B\$1 funciona.} \$A2 fija la columna A: al arrastrar a la derecha sigue leyendo el factor de la fila. B\$1 fija la fila 1: al arrastrar hacia abajo sigue leyendo el factor de la columna. Cada celda multiplica su fila por su columna. Una fórmula, cien resultados.}
\milcpaso{3}{milcMagenta}{\textbf{Cómo evaluar una fórmula ajena.} Pregunta qué se va a copiar y hacia dónde. Mira cada referencia: ¿debe moverse en esa dirección? Si una constante no tiene \$, se rompe. Si un dato de fila tiene \$ en la fila, se repite el mismo valor. Escribe el veredicto con la razón.}
\milcpaso{4}{milcMagenta}{\textbf{Una fórmula, no cien.} Si el cálculo cambia, corriges una celda y vuelves a copiar. Con cien fórmulas escritas a mano, corriges cien y alguna se te escapa. Pensar la referencia antes de copiar es lo que ahorra el trabajo de después.}

\begin{drawbox}{milcMagenta}{Antes de entregar}
\checkbox Hoja del IVA y el descuento con diez productos y las dos fórmulas copiadas. \\[1mm] \checkbox Al cambiar el IVA en E1, la columna C cambia sola. \\[1mm] \checkbox Tabla de multiplicar hecha con una única fórmula en B2. \\[1mm] \checkbox K11 muestra 100 y su fórmula dice \texttt{=\$A11*K\$1}. \\[1mm] \checkbox Captura de la última celda de cada rango. \\[1mm] \checkbox Veredicto sobre la fórmula de mi pareja, con razón.
\end{drawbox}

\begin{softbox}{milcMagenta}{milcGris}{Plantilla de trabajo}
\textbf{IVA.} \textit{=B2*\$E\$1, copiada diez filas.} \\[1mm] \textbf{Descuento.} \textit{=B2*(1-\$E\$2), copiada diez filas.} \\[1mm] \textbf{Tabla de multiplicar.} \textit{=\$A2*B\$1, copiada a cien celdas.} \\[1mm] \textbf{Prueba.} \textit{Cambio el IVA y miro la columna C.} \\[1mm] \textbf{Veredicto.} \textit{La fórmula de mi pareja se rompe o no, porque…}
\end{softbox}

\begin{infoband}{milcMagenta}


\textbf{Lo que va al cuaderno (Actividad 3 · EVALÚA).} \textbf{Título:} «Actividad 3 --- Cien celdas con una fórmula». \textbf{Formato:} la fórmula \texttt{=\$A2*B\$1} explicada en dos líneas (qué fija cada \$), más el veredicto sobre la fórmula de tu pareja con su razón. \textbf{Extensión:} media página. \textbf{Sabes que terminaste cuando} K11 muestra 100, puedes decir qué fija cada \$ de la fórmula, y tu veredicto sobre la fórmula ajena tiene razón escrita.
\end{infoband}

\puente{Lo que entregas hoy: la hoja del IVA y el descuento, la tabla de multiplicar con una única fórmula, y las capturas del final del rango. La prueba de calidad: cambias el IVA en una celda y las diez filas se actualizan solas.}

\milcestacion{milcMostaza}{Producto de la sesión}
\textbf{IVA y descuento con fórmula copiada + tabla de multiplicar con una única fórmula + veredicto}.
\begin{softbox}{milcMostaza}{milcGris}{Criterios de calidad}
(1) El IVA y el descuento se calculan para diez productos con referencias absolutas en las constantes. (2) Al cambiar el IVA en una celda, las diez filas se actualizan solas. (3) La tabla de multiplicar sale de una única fórmula copiada a cien celdas. (4) K11 muestra 100 y su fórmula apunta a A11 y K1. (5) Hay captura de la última celda de cada rango. (6) El veredicto sobre la fórmula de tu pareja tiene razón escrita.
\end{softbox}

\puente{Antes de cerrar, mira lo que hiciste desde cinco lados. Escribir una fórmula que se copia bien es pensar antes de copiar, como la caladora piensa la rejilla antes de la primera puntada.}

\milcestacion{milcVino}{Evaluación liberadora · cinco dimensiones}

\begin{softbox}{milcVino}{milcGris}{Sentido de la evaluación}
Evaluar no es solo revisar una nota. Es reconocer qué aprendí, qué cambió en mí, qué emociones manejé, cómo me relaciono con la ciudad y la comunidad, y qué le dejaré a quien viene detrás.
\end{softbox}

\textbf{Mi reflexión liberadora en cinco dimensiones.} Completa cada frase iniciada:

\small
\begin{tabularx}{\textwidth}{>{\bfseries\RaggedRight\arraybackslash}p{3.4cm}Y>{\RaggedRight\arraybackslash}p{4.6cm}}
\rowcolor{milcVino}\color{white}Dimensión & \color{white}\bfseries Mi respuesta (frame starter) & \color{white}\bfseries Algo que descubrí\\
1. Personal & Lo que más cambió en mí fue \linead{6.5cm} & \linead{4.2cm}\\
2. Emocional & La emoción más fuerte fue \linead{2.5cm} y la regulé \linead{3cm} & \linead{4.2cm}\\
3. Ciudadana & En democracia esto importa porque \linead{6cm} & \linead{4.2cm}\\
4. Local (Cartago) & En mi barrio sería útil para \linead{6.5cm} & \linead{4.2cm}\\
5. Intergeneracional & A mi abuela/abuelo le diría \linead{6.5cm} & \linead{4.2cm}\\
\end{tabularx}
\normalsize

\begin{infoband}{milcVino}
\textbf{Para la próxima sustentación.} Vuelve a esta tabla en seis meses. Verás que las respuestas cambian — no porque lo aprendido cambie, sino porque tú cambias. La evaluación liberadora es una conversación contigo a través del tiempo.
\end{infoband}


\puente{Para terminar, tres ideas sobre lo que se fija y lo que se mueve, contadas con palabras de esta guía: Dussel sobre la tecnología que tiene rostro humano, Epicteto sobre lo que depende de uno, Floridi sobre saber qué preguntarle a los datos.}

\milcestacion{milcGradeDark}{Tres ideas para cerrar · Dussel, un estoico y Floridi}

\begin{softbox}{milcVino}{milcGris}{Enrique Dussel · lente del nosotros}
\textit{No hay liberación sin una tecnología con rostro humano, diseñada desde la historia de la gente que la usa.}

{\footnotesize\itshape\color{milcNegro!70}--- idea tomada de Enrique Dussel, contada con palabras de esta guía. No es una cita textual.}

Una fórmula que se copia bien ahorra horas a quien viene después: al compañero que hereda tu hoja, al tendero que actualiza sus precios. Pensar la referencia antes de copiar es diseñar para otros, no solo para que te dé el resultado.

\textbf{¿Quién va a usar mi hoja después de mí, y le dejé una fórmula que puede corregir?}
\end{softbox}
\linead{\linewidth}\\[1mm]
\linead{\linewidth}

\begin{softbox}{milcMostaza}{milcGris}{Estoicismo · Epicteto · cuidado interior}
\textit{Hay cosas que dependen de nosotros y cosas que no. Saber cuáles son cuáles es el primer trabajo.}

{\footnotesize\itshape\color{milcNegro!70}--- idea tomada de Epicteto, contada con palabras de esta guía. No es una cita textual.}

En una fórmula, el \$ es exactamente eso: decidir qué no va a cambiar aunque copies cien veces, y qué sí. La caladora lo sabe de la rejilla. Tú lo sabes de la constante del IVA.

\textbf{En mi fórmula, ¿qué fijé porque no depende de la fila, y qué dejé libre porque sí?}
\end{softbox}
\linead{\linewidth}\\[1mm]
\linead{\linewidth}

\begin{softbox}{milcTurquesa}{milcGris}{Luciano Floridi · lente de la infoesfera}
\textit{Ganan quienes saben preguntar y responder, y por eso saben qué datos buscar.}

{\footnotesize\itshape\color{milcNegro!70}--- idea tomada de Luciano Floridi, contada con palabras de esta guía. No es una cita textual.}

Antes de copiar, la pregunta es una: ¿qué debe quedarse fijo y qué debe moverse? Quien se la hace escribe una fórmula. Quien no, escribe cien y no sabe cuál falló.

\textbf{¿Me hice la pregunta antes de arrastrar, o arrastré y después miré qué salió?}
\end{softbox}
\linead{\linewidth}\\[1mm]
\linead{\linewidth}

\begin{infoband}{milcNegro}
\textbf{Nota para el docente.} Las tres ideas están contadas con palabras de esta guía a partir del pensamiento de cada autor; no son citas textuales. De 9.º en adelante se trabajan con la obra y su referencia.
\end{infoband}

\milcestacion{milcVino}{Mi autoevaluación · marca una casilla por fila}

{\small
\setlength{\extrarowheight}{2.2pt}
\begin{tabularx}{\textwidth}{>{\RaggedRight\arraybackslash\bfseries}p{3.0cm}YYY}
\rowcolor{milcVino}\color{white}\bfseries Criterio & \color{white}\bfseries Lo logré & \color{white}\bfseries En proceso & \color{white}\bfseries Necesito apoyo\\[.6mm]
Comprensión del tema & \milcopcion{Explico las ideas con mis palabras.} & \milcopcion{Reconozco las ideas, falta precisión.} & \milcopcion{Necesito repasar lo básico.}\\[.8mm]
\rowcolor{milcGris}Pensamiento computacional & \milcopcion{Usé los pilares al armar mi producto.} & \milcopcion{Usé algunos pilares.} & \milcopcion{Necesito releer la estación 2.}\\[.8mm]
Aplicación y producto & \milcopcion{Mi producto cumple los criterios.} & \milcopcion{Está armado, con ajustes pendientes.} & \milcopcion{Aún está en construcción.}\\[.8mm]
\rowcolor{milcGris}Reflexión 5 dimensiones & \milcopcion{Respondí cada una con honestidad.} & \milcopcion{Respondí algunas con detalle.} & \milcopcion{Mis respuestas son muy generales.}\\[.8mm]
Triángulo de pensamiento & \milcopcion{Conecté las 3 voces con mi trabajo.} & \milcopcion{Conecté al menos una.} & \milcopcion{No las relaciono con mi proceso.}\\[.8mm]
\end{tabularx}}

\vspace{3mm}
\textbf{Hoy entendí que:}\\[1mm]
\linead{\linewidth}\\[1mm]
\linead{\linewidth}

\textbf{Mi compromiso para la próxima sesión es:}\\[1mm]
\linead{\linewidth}\\[1mm]
\linead{\linewidth}

\begin{softbox}{milcMostaza}{milcGris}{Cita de despedida · Marco Aurelio}
\textit{``El obstáculo en el camino se vuelve el camino. Lo que estorba al camino se vuelve camino.''}\\[1mm]
Cada error de hoy, bien observado, se convierte en pieza del camino que pisarás mañana.
\end{softbox}

\small
\begin{tabularx}{\textwidth}{>{\bfseries}p{3.6cm}Y}
\rowcolor{milcGradeDark!12}Nombre completo: & \linead{10cm}\\
Fecha: & \linead{4cm} \quad Curso/grupo: \linead{4cm}\\
Mi firma: & \linead{9cm}\\
\end{tabularx}
\normalsize

\begin{infoband}{milcGradeDark}
\textbf{Para la próxima sesión.} Llega con tus dos hojas terminadas. En la sesión 6 vas a combinar varios operadores en una sola fórmula y a aprender en qué orden calcula Excel. Las referencias de hoy son la base; los operadores de mañana se montan sobre ellas.
\end{infoband}

\end{document}
